\section{Experimental Results}
\label{sec:results}

\subsection{Clean validation and held-out test}

Table~\ref{tab:main-results} reports all completed baselines.  The three
traditional rows now represent distinct inputs.  The primary 3,704-D
colour--texture--HOG--BoVW LinearSVC reaches
\PrimaryTraditionalTestTopOnePercent\% test top-1,
\PrimaryTraditionalTestTopFivePercent\% top-5, and
\PrimaryTraditionalTestMacroFOnePercent\% macro F1.  The 2,424-D, 300-tree
handcrafted-only Random Forest without BoVW reaches
\RandomForestTestTopOnePercent\%,
\RandomForestTestTopFivePercent\%, and
\RandomForestTestMacroFOnePercent\%, while the 1,280-D BoVW-only LinearSVC
reaches \BoVWOnlyTestTopOnePercent\%, \BoVWOnlyTestTopFivePercent\%, and
\BoVWOnlyTestMacroFOnePercent\%.  These values are smaller than the learned
representations: scratch ResNet18 reaches 29.88\% top-1, pretrained ResNet18
reaches 67.60\%, and pretrained EfficientNet-B0 reaches
\BestTestTopOnePercent\% top-1, \BestTestTopFivePercent\% top-5, and
\BestTestMacroFOnePercent\% macro F1.  Validation/test rank agreement is
descriptive only; it is not evidence that validation overfitting or prior test
access did not occur.

% Generated from frozen evidence. Do not edit by hand.
\begin{table*}[t]
\centering
\caption{Validation and held-out test results. All values are fractions and are derived from the current hash-bound evidence. The rewritten Stage~2 configuration is frozen going forward; its pre-test locking chronology cannot be independently established retrospectively.}
\label{tab:main-results}
\small
\begin{tabularx}{\textwidth}{@{}Xccccccc@{}}
\toprule
Method & Pretr. & Val top-1 & Val top-5 & Val F1 & Test top-1 & Test top-5 & Test F1 \\
\midrule
Handcrafted + BoVW (LinearSVC) & No & 0.1294 & 0.2702 & 0.115401 & 0.1252 & 0.2678 & 0.111589 \\
Random Forest & No & 0.0860 & 0.2144 & 0.068372 & 0.0822 & 0.2026 & 0.065407 \\
SIFT-BoVW (LinearSVC) & No & 0.0472 & 0.1130 & 0.041271 & 0.0476 & 0.1082 & 0.041405 \\
ResNet18 (scratch) & No & 0.2874 & 0.5274 & 0.279552 & 0.2988 & 0.5364 & 0.292782 \\
ResNet18 (ImageNet pretrained) & Yes & 0.6754 & 0.8646 & 0.673652 & 0.6760 & 0.8624 & 0.672752 \\
EfficientNet-B0 (ImageNet pretrained) & Yes & 0.7368 & 0.9020 & 0.732628 & 0.7274 & 0.9026 & 0.723535 \\
\bottomrule
\end{tabularx}
\end{table*}

% Generated from frozen evidence. Do not edit by hand.
\begin{table*}[t]
\centering
\caption{Held-out macro metrics and measured fit/test timings. Overall accuracy equals top-1 for this single-label task; timing boundaries differ between cached-feature and end-to-end neural evaluation (Section~\ref{sec:setup}).}
\label{tab:test-details}
\small
\begin{tabularx}{\textwidth}{@{}Xccccr@{}}
\toprule
Method & Macro precision & Macro recall & Macro F1 & Train (s) & Test (s) \\
\midrule
Handcrafted + BoVW (LinearSVC) & 0.109007 & 0.125200 & 0.111589 & 1323.2 & 0.190 \\
Random Forest & 0.068343 & 0.082200 & 0.065407 & 870.5 & 1.370 \\
SIFT-BoVW (LinearSVC) & 0.038986 & 0.047600 & 0.041405 & 378.8 & 0.075 \\
ResNet18 (scratch) & 0.305342 & 0.298800 & 0.292782 & 2772.5 & 7.027 \\
ResNet18 (ImageNet pretrained) & 0.691092 & 0.676000 & 0.672752 & 1453.2 & 6.856 \\
EfficientNet-B0 (ImageNet pretrained) & 0.739866 & 0.727400 & 0.723535 & 2421.3 & 8.832 \\
\bottomrule
\end{tabularx}
\end{table*}


\begin{figure*}[t]
  \centering
  \includegraphics[width=0.92\textwidth]{final_model_comparison.png}
  \caption{Held-out comparison of the six complete methods.  The large gap
  between fixed descriptors and fine-tuned pretrained networks dominates the
  smaller within-family differences.}
  \label{fig:model-comparison}
\end{figure*}

The best network is also parameter-efficient: EfficientNet-B0 has 4.65M
trainable parameters, versus 11.43M for ResNet18, yet improves test top-1 by
5.14 points over the selected pretrained ResNet18.  Its canonical training
loop takes 2,421.3 seconds, compared with 1,453.2 seconds for that ResNet18 and
2,772.5 seconds for scratch ResNet18.  These times reflect different epoch
counts and convergence behaviour, not architecture throughput in isolation.
Figure~\ref{fig:model-comparison} visualises the clean-score separation.

\subsection{Validation-only ablations}

Table~\ref{tab:ablations} contains only validation results, so none of these
comparisons consumes additional test information.  The traditional rows show
the selected full four-block validation reference beside colour-only,
texture-only, and HOG-only controls.  The complete 19-configuration grid also
varied the power exponent, block-wise versus global normalisation, feature
subsets, and the three canonical model roles.  It selected block-wise
$\alpha{=}0.3$ with all four blocks; single-block scores must not be interpreted
as a block's contribution in the combined model.

% Generated from frozen validation evidence. Do not edit by hand.
\begin{table*}[t]
\centering
\caption{Validation-only traditional and deep-learning ablations.}
\label{tab:ablations}
\small
\begin{tabularx}{\textwidth}{@{}lXccc@{}}
\toprule
Study membership & Experiment & Val top-1 & Val top-5 & Val macro-F1 \\
\midrule
Feature reference & All blocks + BoVW & 0.1294 & 0.2702 & 0.115401 \\
Feature ablation & Colour only & 0.0608 & 0.1542 & 0.040555 \\
Feature ablation & Texture (LBP) only & 0.0356 & 0.1188 & 0.017688 \\
Feature ablation & HOG only & 0.0292 & 0.0852 & 0.019803 \\
Pretraining & ResNet18 scratch / basic & 0.2874 & 0.5274 & 0.279552 \\
Pretrain / aug. / tuning ref. & ResNet18 pretrained / basic & 0.6616 & 0.8538 & 0.660844 \\
Augmentation & ResNet18 pretrained / strong & 0.6754 & 0.8646 & 0.673652 \\
Frozen head & ResNet18 frozen backbone & 0.5746 & 0.7934 & 0.572458 \\
\bottomrule
\end{tabularx}
\end{table*}


The matched ResNet18 augmentation ablation is the cleanest deep comparison:
strong augmentation improves validation top-1 from 66.16\% to 67.54\%
(+1.38 points) and macro F1 from 66.08\% to 67.37\% (+1.28 points).  The
scratch/basic run reaches 28.74\% and the pretrained/basic run 66.16\%; this
large descriptive gap is consistent with the value of ImageNet transfer, but
is not a one-factor causal estimate because the canonical learning rate,
epoch cap, and stopping policy also differ.  Similarly, the frozen-head
configuration reaches 57.46\% versus 66.16\% for full fine-tuning, but uses a
different optimisation schedule.  It nevertheless shows that the pretrained
representation is already useful before backbone adaptation.

\subsection{Confusions and class-level failures}

The normalised 500-by-500 confusion overview in
Figure~\ref{fig:confusion} is strongly diagonal, but local off-diagonal errors
remain.  The most frequent symmetric pair is Indian Pond Heron versus Squacco
Heron, with six errors (two in one direction and four in the other).  Several
other leading pairs are taxonomically or visually close: Red-eyed Damselfly
versus Common Bluetail, Yellow-faced versus Singing Honeyeater, and
MacGillivray's versus Orange-crowned Warbler.  Figure~\ref{fig:confused-pairs}
preserves both error directions rather than collapsing them into a single
uninterpretable count.

\begin{figure}[t]
  \centering
  \includegraphics[width=0.92\columnwidth]{efficientnet_confusion.png}
  \caption{Row-normalised held-out confusion matrix for EfficientNet-B0
  (rows true, columns predicted).  The overview is complemented by the exact
  integer confused-pair counts in Figure~\ref{fig:confused-pairs}.}
  \label{fig:confusion}
\end{figure}

\begin{figure*}[t]
  \centering
  \includegraphics[width=0.88\textwidth]{top_confused_pairs.png}
  \caption{Ten most frequent bidirectional species confusions for
  EfficientNet-B0.  Colour separates the two directions of each pair.}
  \label{fig:confused-pairs}
\end{figure*}

Class-level recall ranges from 0.1 for totara, Small-leaved Blinks, and New
Zealand jasmine to 1.0 for multiple visually distinctive moth and beetle
classes.  These extremes should be interpreted cautiously because each recall
uses only ten test images.  High-confidence failures are more revealing.  For
example, a matrimony-vine image is predicted as Jerusalem cherry with 0.986
probability, and a Red Junglefowl is predicted as white pasqueflower with
0.932 probability.  Such errors motivate the spatial analysis below.

\subsection{Robustness to degradation}

Table~\ref{tab:robustness} and Figure~\ref{fig:robustness} are generated only
from a complete set of 16 corrupted conditions per model.  EfficientNet-B0 averages
\BestMeanCorruptedTopOnePercent\% top-1 and
\BestMeanCorruptedMacroFOnePercent\% macro F1; its lowest recorded top-1 is
\BestWorstCorruptedTopOnePercent\%, and its lowest recorded macro F1 is
\BestWorstCorruptedMacroFOnePercent\%.  These minima are computed independently
and need not occur under the same corruption.  The traditional curve is not
the primary 3,704-D LinearSVC.  It is the canonical 2,424-D, 300-tree
handcrafted-only Random Forest comparator without BoVW, averaging
\RandomForestMeanCorruptedTopOnePercent\%
top-1 and \RandomForestMeanCorruptedMacroFOnePercent\% macro F1, with worst
recorded top-1 and macro-F1 values of
\RandomForestWorstCorruptedTopOnePercent\% and
\RandomForestWorstCorruptedMacroFOnePercent\%, respectively; these minima are
also computed independently.  The build emits every value in
this paragraph, the table, and the curves from the same complete 64-cell
evidence matrix after the Stage~2 model change.

% Generated from frozen Stage 4 evidence. Do not edit by hand.
\begin{table*}[t]
\centering
\caption{Fixed-protocol robustness summary over 16 corrupted test conditions per model. The traditional comparator is the 2,424-D handcrafted Random Forest, not the primary 3,704-D LinearSVC.}
\label{tab:robustness}
\small
\begin{tabularx}{\textwidth}{@{}Xcccccc@{}}
\toprule
Model & Clean F1 & Mean corrupt. F1 & Worst F1 & Mean $\Delta$ top-1 & Max $\Delta$ top-1 & Rank \\
\midrule
EfficientNet-B0 pretrained & 0.723535 & 0.562260 & 0.191962 & 0.164025 & 0.538000 & 1 \\
ResNet18 pretrained & 0.672752 & 0.478144 & 0.091338 & 0.198675 & 0.587800 & 2 \\
ResNet18 scratch & 0.292782 & 0.216528 & 0.077121 & 0.076512 & 0.211200 & 3 \\
RF (2,424-D handcrafted) & 0.065407 & 0.048932 & 0.014465 & 0.019550 & 0.058600 & 4 \\
\bottomrule
\end{tabularx}
\end{table*}


\begin{figure*}[t]
  \centering
  \includegraphics[width=0.97\textwidth]{robustness_curves.png}
  \caption{Top-1 and macro-F1 trajectories under four deterministic
  corruptions.  Severity zero is the same frozen clean test evaluation for
  every curve.  The traditional curve is the 2,424-D, 300-tree handcrafted
  Random Forest without BoVW, not the primary 3,704-D LinearSVC.  No model is
  retrained or adapted.}
  \label{fig:robustness}
\end{figure*}

The absolute results demonstrate that clean accuracy is not a robustness
guarantee.  The observed pretrained deep models retain higher clean and
corrupted performance in the regenerated evidence, yet the gap between clean
and worst-condition scores remains substantial.  Robustness therefore remains
a first-order failure mode rather than a minor correction to the clean ranking.

\subsection{Grad-CAM and qualitative failure analysis}

Figure~\ref{fig:gradcam} presents predicted-target overlays drawn only from the
pre-registered analysis set.  Correct predictions often concentrate on the
organism: fruit clusters, flower heads, wing markings, or bird bodies.  This
localisation is encouraging but not sufficient evidence of the right
fine-grained cue.  The error cases are more diagnostic.  Both matrimony vine
and Jerusalem cherry maps emphasise red fruit, suggesting that shared colour
and shape dominate subtler botanical distinctions.  The Red Junglefowl map
focuses on a bright textured region and supports a plausible texture shortcut
for the flower prediction.  For the shield-bug/tick error, the predicted map
is tightly centred on the arthropod while the true-target map shifts, showing
that correct localisation can coexist with incorrect classification.

\begin{figure*}[t]
  \centering
  \includegraphics[width=0.32\textwidth]{gradcam_correct_examples.jpg}\hfill
  \includegraphics[width=0.32\textwidth]{gradcam_incorrect_examples.jpg}\hfill
  \includegraphics[width=0.32\textwidth]{gradcam_confused_pairs.jpg}
  \caption{Persisted EfficientNet-B0 Grad-CAM overlays.  Left: four
  high-confidence correct predictions.  Centre: four high-confidence errors.
  Right: representatives of frequent confused pairs.  Captions show true
  $\rightarrow$ predicted common names and softmax
  confidence.  These are qualitative diagnostics, not causal explanations.}
  \label{fig:gradcam}
\end{figure*}

Confused-pair maps likewise centre on the relevant animals, including the two
herons and two damselflies.  Their failure is therefore not simply background
attention; it is consistent with insufficient resolution or feature
specificity for closely related morphology.  Conversely, some overlays span
large rectangular crop regions, reflecting the coarse spatial resolution of
the final convolutional layer.  We consequently treat the maps as spatial
associations rather than evidence of pixel-level causation.
